% Generated draft: requires mathematical and source review.
\documentclass[11pt]{article}
\usepackage[margin=25mm]{geometry}
\usepackage{fontspec}
\setmainfont{Latin Modern Roman}
\usepackage{amsmath,amssymb,mathtools,unicode-math}
\setmathfont{Latin Modern Math}
\usepackage{enumitem,needspace,xurl,hyperref}
\setlength{\parindent}{0pt}
\setlength{\parskip}{5pt}
\setlength{\emergencystretch}{3em}
\newcommand{\N}{\mathbb{N}}
\newcommand{\Z}{\mathbb{Z}}
\newcommand{\Q}{\mathbb{Q}}
\newcommand{\R}{\mathbb{R}}
\newcommand{\C}{\mathbb{C}}
\newcommand{\F}{\mathbb{F}}
\newcommand{\A}{\mathbb{A}}
\newcommand{\PP}{\mathbb{P}}
\newcommand{\E}{\mathbb{E}}
\newcommand{\eps}{\varepsilon}
\newcommand{\dd}{\,\mathrm{d}}
\begin{document}
\section*{UnsolvedMath ID: 30000082}
\textbf{Source problem code:} catalogue entry
\subsection{Definitions and mathematical statement}
Let \(K\) be a finite simplicial complex on a finite vertex set \([n]\): it is a family of subsets closed under taking subsets. Its Stanley--Reisner ideal over a field \(k\) is
\[
I_K=(x_{i_1}\cdots x_{i_r}:\{i_1,\ldots,i_r\}\notin K)
\subset k[x_1,\ldots,x_n].
\]
For a chosen term order over an appropriate infinite field, a generic initial ideal is the stable initial ideal obtained after a generic linear change of coordinates. Algebraic shifting uses generic initial constructions to associate a shifted complex to \(K\). Symmetric and exterior shifting are different constructions; in the symmetric construction a squarefree conversion is part of passing from a generic initial ideal to a complex. One must not simply identify an arbitrary polynomial generic initial ideal with a Stanley--Reisner ideal.

Write \(\Delta_{\mathrm{GIN}}(K)\) for the specific algebraic shifting operation intended in the report, and write \(\operatorname{HVC}(K)\) for its intended homological VC dimension. The source conjecture is the preservation assertion
\[
\forall K\ \forall d\qquad
\operatorname{HVC}(K)=d\ \Longrightarrow\
\operatorname{HVC}(\Delta_{\mathrm{GIN}}(K))=d.
\]
Thus the intended claim is equality of the invariant, not merely an upper or lower bound. The field, coefficient convention for homology, shift variant, term order, and range of \(d\) are not specified in the supplied statement. Most importantly, no definition of HVC dimension is reproduced. The display consequently retains that named invariant rather than replacing it by ordinary VC dimension, a Leray number, or a guessed maximum homology degree. A self-contained definition and precise shifting convention must be recovered from the cited source before this becomes a fully defined theorem statement.

The supplied literature assessment, checked 2026-08-21, says that no proof of this extracted HVC-dimension conjecture was located, and that related homological VC-dimension implications continued to be treated as research questions in the checked literature.
\subsection{Short English statement}
Does the report's generic-initial-ideal algebraic shifting preserve its homological VC dimension exactly? The source record does not define that invariant or fix the version of shifting, so both definitions must be recovered before the conjecture can be fully specified.
\subsection{Sources}
\begin{itemize}
\item Local UnsolvedMath registry, permanent ID 30000082, catalogue entry, \emph{Algebraic Shifting and Homological VC Dimension}.
\item \emph{Combinatorial Commutative Algebra} (2005), Oberwolfach Reports, pp. 1703--1746, DOI \url{https://doi.org/10.4171/owr/2004/32}.
\item Kalai, supplied label \emph{Combinatorial Expectations}: \url{https://gilkalai.wordpress.com/wp-content/uploads/2019/03/ow04.pdf}.
\item Supplied literature triage checked 2026-08-21; the statement-verification field reports comparison with the official report, without reproducing the missing definitions.
\end{itemize}
\end{document}
